\documentclass[11pt,a4paper]{article}

% ============================================================
% ENCODAGE & POLICES
% ============================================================
\usepackage[T1]{fontenc}
\usepackage[utf8]{inputenc}
\usepackage{mathpazo}          % Palatino (elegant, lisible)
\usepackage[scaled=0.88]{helvet} % Helvetica pour sans-serif
\usepackage{microtype}         % Optimisation typographique

% ============================================================
% MISE EN PAGE
% ============================================================
\usepackage[
  top=2.8cm, bottom=2.8cm,
  left=2.5cm, right=2.5cm,
  headheight=14pt
]{geometry}
\usepackage{parskip}           % Espacement entre paragraphes
\usepackage{setspace}
\setstretch{1.15}              % Interligne agreable

% ============================================================
% COULEURS (palette sobre et moderne)
% ============================================================
\usepackage[dvipsnames,svgnames]{xcolor}
\definecolor{deepblue}{HTML}{1D3557}    % Bleu profond (titres)
\definecolor{accent}{HTML}{457B9D}      % Bleu moyen (sous-titres)
\definecolor{lightblue}{HTML}{A8DADC}   % Bleu clair (fond boites)
\definecolor{warmgray}{HTML}{F1FAEE}    % Gris chaud (fond def)
\definecolor{alertred}{HTML}{E63946}    % Rouge alerte
\definecolor{softgreen}{HTML}{2D6A4F}   % Vert methode
\definecolor{lightgreen}{HTML}{D8F3DC}  % Vert clair fond
\definecolor{lightyellow}{HTML}{FFF8E7} % Jaune clair fond
\definecolor{lightred}{HTML}{FFE8E8}    % Rouge clair fond
\definecolor{textmain}{HTML}{2B2D42}    % Texte principal
\definecolor{textlight}{HTML}{6C757D}   % Texte secondaire

% ============================================================
% EN-TETE & PIED DE PAGE
% ============================================================
\usepackage{fancyhdr}
\pagestyle{fancy}
\fancyhf{}
\renewcommand{\headrulewidth}{0.4pt}
\renewcommand{\headrule}{\hbox to\headwidth{\color{lightblue}\leaders\hrule height \headrulewidth\hfill}}
\fancyhead[L]{\small\sffamily\color{textlight}Chapitre 6}
\fancyhead[R]{\small\sffamily\color{textlight}Tables statistiques : lois continues}
\fancyfoot[C]{\small\sffamily\color{textlight}\thepage}

% ============================================================
% TITRES (design epure)
% ============================================================
\usepackage{titlesec}

\titleformat{\section}
  {\normalfont\Large\bfseries\sffamily\color{deepblue}}
  {\thesection.}{0.6em}{}
  [\vspace{-0.3em}{\color{lightblue}\rule{\textwidth}{1.5pt}}]

\titleformat{\subsection}
  {\normalfont\large\bfseries\sffamily\color{accent}}
  {\thesubsection.}{0.5em}{}

\titleformat{\subsubsection}
  {\normalfont\normalsize\bfseries\color{textmain}}
  {\thesubsubsection.}{0.5em}{}

\titlespacing*{\section}{0pt}{2em}{1em}
\titlespacing*{\subsection}{0pt}{1.5em}{0.6em}
\titlespacing*{\subsubsection}{0pt}{1em}{0.4em}

% ============================================================
% BOITES (tcolorbox -- design minimal)
% ============================================================
\usepackage[most]{tcolorbox}

% Definition
\newtcolorbox{defbox}[1][]{%
  enhanced, breakable,
  colback=warmgray,
  colframe=deepblue,
  boxrule=0pt, leftrule=3pt,
  arc=0pt, outer arc=0pt,
  fonttitle=\bfseries\sffamily\color{deepblue},
  title={#1},
  top=8pt, bottom=8pt, left=12pt, right=12pt,
  before skip=10pt, after skip=10pt,
  toptitle=4pt, bottomtitle=4pt,
}

% Theoreme
\newtcolorbox{thmbox}[1][]{%
  enhanced, breakable,
  colback=white,
  colframe=accent,
  boxrule=0.6pt,
  arc=3pt, outer arc=3pt,
  fonttitle=\bfseries\sffamily\color{accent},
  title={#1},
  top=8pt, bottom=8pt, left=12pt, right=12pt,
  before skip=10pt, after skip=10pt,
  toptitle=4pt, bottomtitle=4pt,
}

% Methode
\newtcolorbox{methbox}[1][]{%
  enhanced, breakable,
  colback=lightgreen,
  colframe=softgreen,
  boxrule=0pt, leftrule=3pt,
  arc=0pt, outer arc=0pt,
  fonttitle=\bfseries\sffamily\color{softgreen},
  title={#1},
  top=8pt, bottom=8pt, left=12pt, right=12pt,
  before skip=10pt, after skip=10pt,
  toptitle=4pt, bottomtitle=4pt,
}

% Attention / piege
\newtcolorbox{alertbox}[1][]{%
  enhanced, breakable,
  colback=lightred,
  colframe=alertred,
  boxrule=0pt, leftrule=3pt,
  arc=0pt, outer arc=0pt,
  fonttitle=\bfseries\sffamily\color{alertred},
  title={#1},
  top=8pt, bottom=8pt, left=12pt, right=12pt,
  before skip=10pt, after skip=10pt,
  toptitle=4pt, bottomtitle=4pt,
}

% Exemple
\newtcolorbox{exbox}[1][]{%
  enhanced, breakable,
  colback=lightyellow,
  colframe=orange!60!yellow,
  boxrule=0pt, leftrule=3pt,
  arc=0pt, outer arc=0pt,
  fonttitle=\bfseries\sffamily\color{orange!70!black},
  title={#1},
  top=8pt, bottom=8pt, left=12pt, right=12pt,
  before skip=10pt, after skip=10pt,
  toptitle=4pt, bottomtitle=4pt,
}

% ============================================================
% MATHS
% ============================================================
\usepackage{amsmath,amssymb,amsthm}
\usepackage{mathtools}

% ============================================================
% TABLEAUX
% ============================================================
\usepackage{booktabs}
\usepackage{array}
\usepackage{tabularx}
\usepackage{multirow}

% ============================================================
% LISTES
% ============================================================
\usepackage{enumitem}
\setlist[itemize]{leftmargin=1.2em, label={\color{accent}\textbullet}, itemsep=3pt}
\setlist[enumerate]{leftmargin=1.5em, itemsep=3pt}

% ============================================================
% LIENS
% ============================================================
\usepackage{hyperref}
\hypersetup{
  colorlinks=true,
  linkcolor=deepblue,
  urlcolor=accent,
  citecolor=accent,
  pdfauthor={},
  pdftitle={Chapitre 6 -- Tables statistiques : lois continues},
}

% ============================================================
% COMMANDES PERSO
% ============================================================
\newcommand{\N}{\mathcal{N}}
\newcommand{\E}{\mathbb{E}}
\newcommand{\Var}{\mathrm{Var}}
\newcommand{\Prob}{\mathbb{P}}
\newcommand{\R}{\mathbb{R}}
\renewcommand{\P}{\Prob}

% ============================================================
\begin{document}

% ============================================================
% PAGE DE TITRE
% ============================================================
\thispagestyle{empty}
\begin{center}
\vspace*{4cm}

{\fontsize{72}{80}\selectfont\sffamily\bfseries\color{deepblue}6}

\vspace{0.5cm}
{\color{lightblue}\rule{8cm}{2pt}}
\vspace{0.6cm}

{\LARGE\sffamily\bfseries\color{deepblue}Tables statistiques :}\\[0.3cm]
{\LARGE\sffamily\bfseries\color{deepblue}lois continues}

\vspace{1.5cm}

{\large\sffamily\color{textlight}%
Loi Normale\quad$\cdot$\quad Khi-deux\quad$\cdot$\quad Student\quad$\cdot$\quad Fisher-Snedecor}

\vspace{2.5cm}

{\sffamily\color{textlight}Cours complet -- L2 \'Economie\\[0.3cm]
Statistiques -- Semestre 2}

\vfill

{\small\sffamily\color{textlight}\today}
\end{center}

\newpage

% ============================================================
% TABLE DES MATI\`ERES
% ============================================================
\thispagestyle{empty}
{\hypersetup{linkcolor=textmain}
\tableofcontents}
\newpage

% ============================================================
% INTRODUCTION
% ============================================================
\section*{Introduction}
\addcontentsline{toc}{section}{Introduction}

Ce chapitre porte sur la \textbf{construction des lois continues} et la \textbf{lecture de leurs tables statistiques}. Les quatre lois \'etudi\'ees sont : la loi Normale, la loi du Khi-deux, la loi de Student et la loi de Fisher-Snedecor.

L'int\'er\^et de ce chapitre est double :
\begin{enumerate}
\item La construction de lois statistiques est utile pour la d\'etermination de \textbf{lois d'\'echantillonnage} (\textit{cf.} Chapitre~8).
\item Les lectures de tables permettent de trouver des probabilit\'es, de d\'eterminer des \textbf{intervalles de confiance} et des \textbf{intervalles d'acceptation} dans le cadre de la th\'eorie des tests (\textit{cf.} Chapitre~11).
\end{enumerate}

\begin{thmbox}[Fiche de synth\`ese -- Constructions de lois]
\textbf{Khi-deux :}\quad Soient $U_1,\ldots,U_n$ i.i.d. $\sim \N(0;1)$. Alors :
$$K = \sum_{i=1}^{n} U_i^2 \;\sim\; \chi^2(n).$$

\textbf{Student :}\quad Soient $U \sim \N(0;1)$ et $K \sim \chi^2(\nu)$ ind\'ependantes. Alors :
$$T = \frac{U}{\sqrt{K/\nu}} \;\sim\; t(\nu).$$

\textbf{Fisher-Snedecor :}\quad Soient $K_1 \sim \chi^2(\nu_1)$ et $K_2 \sim \chi^2(\nu_2)$ ind\'ependantes. Alors :
$$F = \frac{K_1/\nu_1}{K_2/\nu_2} \;\sim\; F(\nu_1;\nu_2).$$
\end{thmbox}

% ============================================================
% I. LOI NORMALE
% ============================================================
\section{La loi Normale}

\subsection{D\'efinition et param\`etres}

\begin{defbox}[D\'efinition -- Loi normale]
Une variable al\'eatoire $X$ suit une \textbf{loi normale} de param\`etres $m$ (esp\'erance) et $\sigma$ (\'ecart-type), not\'ee $X \sim \N(m;\sigma)$, si sa densit\'e de probabilit\'e est :
$$f_X(x) = \frac{1}{\sigma\sqrt{2\pi}}\,\exp\!\left(-\frac{1}{2}\left(\frac{x-m}{\sigma}\right)^{\!2}\right), \qquad x \in \R.$$
On a alors $\E(X)=m$, $\Var(X)=\sigma^2$, et la courbe de $f_X$ est la \og{}courbe en cloche\fg{} sym\'etrique autour de $m$.
\end{defbox}

\subsection{Centrage et r\'eduction}

Pour utiliser une \textbf{unique table} quelle que soit la loi $\N(m;\sigma)$, on effectue le changement de variable :

\begin{defbox}[Changement de variable fondamental]
$$U = \frac{X - m}{\sigma} \qquad\Longrightarrow\qquad U \sim \N(0;1).$$
La variable $U$ est appel\'ee \textbf{variable normale centr\'ee r\'eduite}. Sa densit\'e est :
$$f_U(u) = \frac{1}{\sqrt{2\pi}}\,\exp\!\left(-\frac{u^2}{2}\right).$$
Et on a la relation : $f_X(x) = \dfrac{1}{\sigma}\,f_U\!\left(\dfrac{x-m}{\sigma}\right)$.
\end{defbox}

\begin{alertbox}[Attention]
La notation $\N(m;\sigma)$ utilise l'\textbf{\'ecart-type} $\sigma$, pas la variance $\sigma^2$. V\'erifier syst\'ematiquement ce que donne l'\'enonc\'e.
\end{alertbox}

\subsection{Table 1 -- Densit\'e de la loi $\N(0;1)$}

La Table~1 donne directement $f_U(u)$ pour $u \geq 0$. Gr\^ace \`a la propri\'et\'e de sym\'etrie $f_U(-u) = f_U(u)$, une seule colonne suffit.

\begin{exbox}[Exemple de lecture]
\textbf{Donn\'ee :} $X \sim \N(1{,}75\;;\;0{,}1)$. Calculer $\P(X = 2)$, i.e. $f_X(2)$.

\textbf{M\'ethode :} On centre et r\'eduit : $u = \dfrac{2-1{,}75}{0{,}1} = 2{,}5$.

On lit dans la Table~1 : $f_U(2{,}5) = 0{,}0175$.

Donc : $f_X(2) = \dfrac{1}{\sigma}\,f_U(u) = \dfrac{1}{0{,}1}\times 0{,}0175 = 0{,}175.$
\end{exbox}

\subsection{Table 2 -- Fonction de r\'epartition $F_U(u) = \P(U \leq u)$}

\begin{defbox}[Fonction de r\'epartition]
$$F_U(u) = \P(U \leq u) = \int_{-\infty}^{u} f_U(t)\,dt.$$
La table donne $F_U(u)$ pour $u \geq 0$ avec deux d\'ecimales en ligne et la troisi\`eme en colonne.
\end{defbox}

\subsubsection{Propri\'et\'es fondamentales}

\begin{thmbox}[Propri\'et\'es de $F_U$]
\begin{enumerate}
\item \textbf{Sym\'etrie :}\quad $F_U(-u) = 1 - F_U(u)$\qquad (\textit{ne pas confondre avec} $f(-u) = f(u)$).
\item \textbf{Compl\'ementaire :}\quad $\P(U > u) = 1 - F_U(u)$.
\item \textbf{Intervalle :}\quad $\P(a < U \leq b) = F_U(b) - F_U(a)$.
\item \textbf{Continuit\'e :}\quad $\P(a < U < b) = \P(a \leq U \leq b) = F_U(b) - F_U(a)$.
\end{enumerate}
\end{thmbox}

\begin{exbox}[Exemples de lecture]
\begin{enumerate}
\item $F_U(0{,}22) = \P(U \leq 0{,}22) = 0{,}58706$.
\item $F_U(-0{,}32) = 1 - F_U(0{,}32) = 1 - 0{,}62552 = 0{,}37448$.
\item $F_U(-1{,}5) = 1 - F_U(1{,}5) = 1 - 0{,}93319 = 0{,}06681$.
\item $\P(-1{,}48 < U < 1{,}80) = F_U(1{,}80) - F_U(-1{,}48) = 0{,}9641 - (1 - 0{,}9306) = 0{,}8947$.
\end{enumerate}
\end{exbox}

\begin{exbox}[Exemple avec changement de variable]
$X \sim \N(1;6)$. Calculer $\P(X \leq 4)$.

Centrage-r\'eduction : $u = \dfrac{4-1}{6} = 0{,}5$.

$\P(X \leq 4) = \P\!\left(U \leq 0{,}5\right) = F_U(0{,}5) = 0{,}6915.$
\end{exbox}

\subsection{Table 3 -- Fractiles (quantiles) de la loi $\N(0;1)$}

La Table~3 est la \textbf{table inverse} : on entre avec une probabilit\'e et on obtient le fractile $u$.

\begin{defbox}[Notations des fractiles]
On note $u_{1-P}$ (ou $U_{1-P}$) le fractile tel que :
$$\P(U > u_{1-P}) = P \qquad\Longleftrightarrow\qquad \P(U \leq u_{1-P}) = 1-P = Q.$$
L'indice fait r\'ef\'erence \`a la valeur de la fonction de r\'epartition $Q$, \textbf{pas} \`a $P$.
\end{defbox}

\subsubsection{Lecture unilatérale}

\begin{exbox}[Exemples]
\begin{itemize}
\item $\P(U > u) = 0{,}05$ : on lit $u_{0{,}95} = 1{,}6449$.
\item $\P(U > u) = 0{,}025$ : on lit $u_{0{,}975} = 1{,}96$.
\item Pour $Q < 0{,}5$ (fractile n\'egatif) : $u_Q = -u_{1-Q}$.\\ Exemple : $u_{0{,}05} = -u_{0{,}95} = -1{,}6449$.
\end{itemize}
\end{exbox}

\subsubsection{Lecture bilat\'erale sym\'etrique}

Pour un intervalle $\P(-u < U < u) = 1 - P$, on r\'epartit $P/2$ de chaque c\^ot\'e :

\begin{methbox}[R\`egle de lecture bilat\'erale]
\textbf{Lire la Table~3 \`a $P/2$ au lieu de $P$.}
\begin{itemize}
\item $P = 5\%$ : lire \`a $P/2 = 0{,}025$ $\Rightarrow$ $u = u_{0{,}975} = 1{,}96$.\\ Bornes : $-1{,}96$ et $+1{,}96$.
\item $P = 10\%$ : lire \`a $P/2 = 0{,}05$ $\Rightarrow$ $u = u_{0{,}95} = 1{,}6449$.\\ Bornes : $-1{,}6449$ et $+1{,}6449$.
\item $P = 1\%$ : lire \`a $P/2 = 0{,}005$ $\Rightarrow$ $u = u_{0{,}995} = 2{,}5758$.\\ Bornes : $-2{,}5758$ et $+2{,}5758$.
\end{itemize}
\end{methbox}

\subsection{Th\'eor\`emes d'additivit\'e}

\begin{thmbox}[Th\'eor\`eme d'additivit\'e des lois normales]
\textbf{Hypoth\`eses :} $X_i \sim \N(m_i;\sigma_i)$, $i=1,\ldots,n$, mutuellement ind\'ependantes.

\textbf{R\'esultat :}
$$\sum_{i=1}^{n} X_i \;\sim\; \N\!\left(\sum_{i=1}^{n} m_i\;;\;\sqrt{\sum_{i=1}^{n} \sigma_i^2}\right).$$
\end{thmbox}

\begin{alertbox}[Rappel]
L'\'ecart-type de la somme est $\sqrt{\sum \sigma_i^2}$, \textbf{pas} $\sum \sigma_i$. Les variances s'additionnent, pas les \'ecarts-types.
\end{alertbox}

% ============================================================
% II. LOI DU KHI-DEUX
% ============================================================
\newpage
\section{La loi du Khi-deux}

\subsection{Construction}

\begin{defbox}[D\'efinition -- Loi du Khi-deux]
Soient $U_1,\ldots,U_n$ des variables al\'eatoires ind\'ependantes suivant chacune $\N(0;1)$. La variable al\'eatoire :
$$K = \sum_{i=1}^{n} U_i^2 \;\sim\; \chi^2(n)$$
suit une \textbf{loi du Khi-deux \`a $n$ degr\'es de libert\'e}.

Le \textbf{degr\'e de libert\'e} $\nu = n$ indique le nombre de $U_i^2$ somm\'ees. Par convention, on note le degr\'e de libert\'e $\nu$.
\end{defbox}

\textbf{Contexte d'utilisation :} la loi du Khi-deux intervient naturellement pour les indicateurs de dispersion (variance). En \'economie, elle permet de tester si la variabilit\'e d'un processus de production reste dans les limites du cahier des charges.

\subsection{Lecture de la Table 4}

\begin{defbox}[Notation des fractiles du Khi-deux]
La Table~4 donne le fractile $\chi^2_{1-P}(\nu)$ tel que :
$$\P\!\left(K > \chi^2_{1-P}(\nu)\right) = P.$$
On lit \`a la colonne $P$ et \`a la ligne $\nu$.
\end{defbox}

\subsubsection{Lecture unilatérale}

\begin{exbox}[Exemples]
\begin{itemize}
\item $\chi^2_{0{,}975}(18) = 31{,}53$ : lire colonne $P = 0{,}025$, ligne $\nu = 18$.
\item $\chi^2_{0{,}10}(9) = 4{,}168$ : lire colonne $P = 0{,}9$, ligne $\nu = 9$.
\item $\chi^2_{0{,}99}(30) = 50{,}892$ : lire colonne $P = 0{,}01$, ligne $\nu = 30$.
\end{itemize}
\end{exbox}

\subsubsection{Lecture bilat\'erale sym\'etrique}

Pour un intervalle $\P\!\left(\chi^2_{P/2} < K \leq \chi^2_{1-P/2}\right) = 1-P$, on r\'epartit $P/2$ de chaque c\^ot\'e.

\begin{exbox}[Exemples bilat\'eraux]
\begin{itemize}
\item $P=5\%$, $\nu=18$ : lire \`a $1-P/2 = 0{,}975$ et $P/2 = 0{,}025$.\\
$\chi^2_{0{,}025}(18) = 8{,}23$ \quad et \quad $\chi^2_{0{,}975}(18) = 31{,}53$.
\item $P=10\%$, $\nu=1$ : lire \`a $0{,}95$ et $0{,}05$.\\
$\chi^2_{0{,}05}(1) = 0{,}0039$ \quad et \quad $\chi^2_{0{,}95}(1) = 3{,}841$.
\end{itemize}
\end{exbox}

\subsection{Th\'eor\`eme d'additivit\'e}

\begin{thmbox}[Additivit\'e des lois de Khi-deux]
\textbf{Hypoth\`eses :} $K_i \sim \chi^2(\nu_i)$, $i=1,\ldots,n$, mutuellement ind\'ependantes.

\textbf{R\'esultat :}
$$\sum_{i=1}^{n} K_i \;\sim\; \chi^2\!\left(\sum_{i=1}^{n} \nu_i\right).$$
Les degr\'es de libert\'e s'additionnent.
\end{thmbox}

\subsection{Convergence vers la loi normale ($\nu > 30$)}

Lorsque $\nu > 30$, la loi du Khi-deux converge vers la loi normale. Plus pr\'ecis\'ement :
$$\sqrt{2K} - \sqrt{2\nu-1} \;\overset{\text{approx.}}{\sim}\; \N(0;1).$$

On peut alors approcher le fractile :
$$\chi^2_{1-P}(\nu) \;\approx\; \frac{\left(u_{1-P} + \sqrt{2\nu-1}\right)^2}{2}.$$

\begin{exbox}[Exemple]
$\chi^2_{0{,}90}(70)$ (probabilit\'e $P=10\%$ d'\^etre d\'epass\'ee) :
$$\chi^2_{0{,}90}(70) \approx \frac{(1{,}2816 + \sqrt{139})^2}{2} = \frac{(1{,}2816 + 11{,}79)^2}{2} \approx 85{,}43.$$
\end{exbox}

% ============================================================
% III. LOI DE STUDENT
% ============================================================
\newpage
\section{La loi de Student}

\subsection{Construction}

\begin{defbox}[D\'efinition -- Loi de Student]
Soient $U \sim \N(0;1)$ et $K \sim \chi^2(\nu)$ \textbf{ind\'ependantes}. La variable al\'eatoire :
$$T = \frac{U}{\sqrt{K/\nu}} \;\sim\; t(\nu)$$
suit une \textbf{loi de Student \`a $\nu$ degr\'es de libert\'e}.
\end{defbox}

\textbf{Propri\'et\'es cl\'es :}
\begin{itemize}
\item La densit\'e de Student est \textbf{sym\'etrique en 0}, comme la loi normale.
\item Queues plus \'epaisses que la normale (plus de valeurs extr\^emes).
\item Quand $\nu \to \infty$, $t(\nu) \to \N(0;1)$.
\end{itemize}

\subsection{Lecture de la Table 5}

\begin{defbox}[Convention de la Table de Student]
La Table~5 est directement pr\'evue pour des \textbf{intervalles bilat\'eraux sym\'etriques}. Elle donne le fractile $T_{1-P}(\nu)$ tel que :
$$\P\!\left(|T| > T_{1-P}(\nu)\right) = P.$$
Autrement dit, $\P\!\left(-T_{1-P} < T < T_{1-P}\right) = 1-P$.

Par sym\'etrie : $T_P(\nu) = -T_{1-P}(\nu)$.
\end{defbox}

\begin{alertbox}[Diff\'erence cl\'e avec les autres tables]
Contrairement aux tables de la loi normale et du Khi-deux, la Table de Student est con\c{c}ue pour une lecture bilat\'erale directe. Pour un intervalle \textbf{unilat\'eral} \`a risque $P$, il faut lire \`a la colonne $2P$ (on double $P$).
\end{alertbox}

\begin{exbox}[Exemples de lecture]
\begin{itemize}
\item $T_{0{,}80}(29) = 1{,}3114$ (bilat\'eral) : lire colonne $P=0{,}2$, ligne $\nu=29$.
\item $T_{0{,}70}(2)$ (bilat\'eral) : lire colonne $P = 0{,}30$, ligne $\nu=2$ $\Rightarrow$ $1{,}3862$.\\
Par sym\'etrie : $T_{0{,}30}(2) = -1{,}3862$.
\item Unilat\'eral \`a $P=5\%$ avec $\nu=10$ : lire colonne $2P = 0{,}10$ $\Rightarrow$ $T_{0{,}90}(10)$.
\end{itemize}
\end{exbox}

\subsection{Convergence vers la loi normale ($\nu > 30$)}

\begin{thmbox}[Convergence]
Pour $\nu > 30$, la loi de Student converge vers la loi normale centr\'ee r\'eduite :
$$T \;\overset{\nu > 30}{\approx}\; U \sim \N(0;1).$$
Les fractiles de Student sont alors remplac\'es par ceux de la loi normale :
\begin{itemize}
\item $T_{0{,}95}(\nu) \approx u_{0{,}95} = 1{,}96$ (bilat\'eral sym\'etrique, $P=5\%$).
\item $T_{0{,}90}(\nu) \approx u_{0{,}90} = 1{,}6449$ (bilat\'eral sym\'etrique, $P=10\%$).
\item $T_{0{,}95}(120) = u_{0{,}95} = 1{,}2886$ (le fractile exact est tr\`es proche de la normale).
\end{itemize}
\end{thmbox}

% ============================================================
% IV. LOI DE FISHER-SNEDECOR
% ============================================================
\newpage
\section{La loi de Fisher-Snedecor}

\subsection{Construction}

\begin{defbox}[D\'efinition -- Loi de Fisher-Snedecor]
Soient $K_1 \sim \chi^2(\nu_1)$ et $K_2 \sim \chi^2(\nu_2)$ \textbf{ind\'ependantes}. La variable al\'eatoire :
$$F = \frac{K_1/\nu_1}{K_2/\nu_2} \;\sim\; F(\nu_1;\nu_2)$$
suit une \textbf{loi de Fisher-Snedecor} \`a $\nu_1$ et $\nu_2$ degr\'es de libert\'e.
\end{defbox}

\textbf{Propri\'et\'es :}
\begin{itemize}
\item $F$ est \`a valeurs \textbf{positives} (rapport de deux quantit\'es positives).
\item La loi n'est \textbf{pas sym\'etrique} : l'ordre $(\nu_1;\nu_2)$ compte.
\end{itemize}

\subsection{Lecture de la Table 6}

\begin{defbox}[Convention de la Table de Fisher]
La Table~6 donne le fractile $F_{1-P}(\nu_1;\nu_2)$ tel que :
$$\P\!\left(F > F_{1-P}(\nu_1;\nu_2)\right) = P = 5\%.$$
On lit $\nu_1$ en \textbf{colonne} et $\nu_2$ en \textbf{ligne}.
\end{defbox}

\begin{exbox}[Exemples unilatéraux]
\begin{itemize}
\item $F_{0{,}95}(3;5) \Leftrightarrow \P(F > F_{0{,}95}(3;5)) = 5\%$ : lire colonne $\nu_1=3$, ligne $\nu_2=5$.\\
On trouve $F_{0{,}95}(3;5) = 5{,}41$.
\item $F_{0{,}95}(\infty;100) = 1{,}30$ : colonne $\infty$, ligne $100$.
\end{itemize}
\end{exbox}

\subsection{Propri\'et\'e fondamentale : inversion des degr\'es de libert\'e}

\begin{thmbox}[Formule des queues gauches]
Pour un risque $P < 50\%$ :
$$\boxed{F_P(\nu_1;\nu_2) = \frac{1}{F_{1-P}(\nu_2;\nu_1)}.}$$
Cette formule permet de trouver les fractiles des queues \textbf{gauches} \`a partir d'une table qui ne donne que les queues droites.
\end{thmbox}

\begin{exbox}[Exemples]
\begin{itemize}
\item $F_{0{,}05}(5;3) = \dfrac{1}{F_{0{,}95}(3;5)} = \dfrac{1}{5{,}41} \approx 0{,}18$.\\ Et : $F_{0{,}95}(5;3) = 9{,}01$.
\item $F_{0{,}01}(50;100) = \dfrac{1}{F_{0{,}99}(100;50)} = \dfrac{1}{1{,}82} \approx 0{,}55$.
\end{itemize}
\end{exbox}

\subsection{Lien Student -- Fisher}

\begin{thmbox}[Propri\'et\'e]
Si $T \sim t(\nu)$, alors :
$$T^2 \;\sim\; F(1;\nu).$$
\end{thmbox}

% ============================================================
% V. FICHE DE SYNTHESE
% ============================================================
\newpage
\section{Fiche de synth\`ese}

\begin{center}
\renewcommand{\arraystretch}{1.6}
\sffamily\small
\begin{tabular}{>{\bfseries}p{2.5cm} p{3.8cm} p{2.5cm} p{4.2cm}}
\toprule
\textbf{Loi} & \textbf{Construction} & \textbf{Param\`etre} & \textbf{Lecture de table} \\
\midrule
Normale $\N(0;1)$ &
Centrage-r\'eduction : $U = \frac{X-m}{\sigma}$ &
$m$, $\sigma$ &
Table 2 : $F_U(u)$\newline Table 3 : fractiles $u_{1-P}$ \\
\addlinespace
Khi-deux $\chi^2(\nu)$ &
$K = \sum U_i^2$ &
$\nu$ (ddl) &
Table 4 : $\P(K > \chi^2_{1-P}) = P$\newline Colonne $P$, ligne $\nu$ \\
\addlinespace
Student $t(\nu)$ &
$T = \frac{U}{\sqrt{K/\nu}}$ &
$\nu$ (ddl) &
Table 5 : bilat\'eral direct\newline $\P(|T| > T_{1-P}) = P$ \\
\addlinespace
Fisher $F(\nu_1;\nu_2)$ &
$F = \frac{K_1/\nu_1}{K_2/\nu_2}$ &
$\nu_1$, $\nu_2$ (ddl) &
Table 6 : $\nu_1$ en colonne, $\nu_2$ en ligne\newline Queue gauche : $F_P = 1/F_{1-P}$ \\
\bottomrule
\end{tabular}
\end{center}

\begin{methbox}[Fractiles \`a m\'emoriser]
\renewcommand{\arraystretch}{1.3}
\begin{center}
\begin{tabular}{lcccc}
\toprule
$P$ (bilat\'eral) & $10\%$ & $5\%$ & $2\%$ & $1\%$ \\
\midrule
$u_{1-P/2}$ & $1{,}6449$ & $1{,}96$ & $2{,}3263$ & $2{,}5758$ \\
\bottomrule
\end{tabular}
\end{center}
\end{methbox}

% ============================================================
% VI. METHODES INDISPENSABLES
% ============================================================
\newpage
\section{M\'ethodes indispensables au TD / partiel}

\begin{methbox}[M\'ethode 1 -- Calculer $\P(a < X < b)$ avec $X \sim \N(m;\sigma)$]
\begin{enumerate}
\item Centrer-r\'eduire les bornes : $u_a = \dfrac{a-m}{\sigma}$, $u_b = \dfrac{b-m}{\sigma}$.
\item Lire dans la Table~2 : $F_U(u_b)$ et $F_U(u_a)$.
\item Si $u < 0$ : $F_U(u) = 1 - F_U(-u)$.
\item R\'esultat : $\P(a < X < b) = F_U(u_b) - F_U(u_a)$.
\end{enumerate}
\end{methbox}

\begin{methbox}[M\'ethode 2 -- Trouver un fractile (table inverse)]
\begin{enumerate}
\item Identifier la probabilit\'e donn\'ee et le type d'intervalle (unilat\'eral/bilat\'eral).
\item Si bilat\'eral : diviser $P$ par $2$ avant de lire.
\item Lire le fractile dans la Table~3 (normale), Table~4 (khi-deux) ou Table~5 (Student).
\item Si fractile n\'egatif (normale) : $u_Q = -u_{1-Q}$.
\end{enumerate}
\end{methbox}

\begin{methbox}[M\'ethode 3 -- Construire une loi \`a partir de l'\'enonc\'e]
\begin{enumerate}
\item Identifier les variables al\'eatoires en jeu et leurs lois.
\item V\'erifier l'ind\'ependance (hypoth\`ese toujours n\'ecessaire).
\item Reconna\^itre la structure :
    \begin{itemize}
    \item Somme de $\N(0;1)^2$ $\Rightarrow$ Khi-deux.
    \item Rapport $\N(0;1)/\sqrt{\chi^2/\nu}$ $\Rightarrow$ Student.
    \item Rapport de deux $\chi^2$ normalis\'es $\Rightarrow$ Fisher.
    \end{itemize}
\item D\'eterminer les degr\'es de libert\'e.
\end{enumerate}
\end{methbox}

\begin{methbox}[M\'ethode 4 -- Convergence ($\nu > 30$)]
Quand le degr\'e de libert\'e d\'epasse 30 :
\begin{itemize}
\item \textbf{Student :} remplacer $T_{1-P}(\nu)$ par $u_{1-P}$ (fractile de la normale).
\item \textbf{Khi-deux :} utiliser $\chi^2_{1-P}(\nu) \approx \dfrac{(u_{1-P} + \sqrt{2\nu-1})^2}{2}$.
\end{itemize}
\end{methbox}

% ============================================================
% VII. EXERCICES CORRIGES
% ============================================================
\newpage
\section{Exercices corrig\'es}

% --- Exercice 1 ---
\begin{thmbox}[Exercice 1 -- Lecture de table de la loi normale]
Soit $X \sim \N(100;15)$.
\begin{enumerate}[label=\alph*)]
\item Calculer $\P(X \leq 120)$.
\item Calculer $\P(85 < X < 130)$.
\item Trouver $x$ tel que $\P(X > x) = 0{,}05$.
\end{enumerate}
\end{thmbox}

\textbf{Correction :}

\textbf{a)} Centrage-r\'eduction : $u = \dfrac{120-100}{15} = \dfrac{20}{15} = 1{,}33$.

$\P(X \leq 120) = F_U(1{,}33) = 0{,}90824.$

\medskip
\textbf{b)} On centre : $u_1 = \dfrac{85-100}{15} = -1$ et $u_2 = \dfrac{130-100}{15} = 2$.

$\P(85 < X < 130) = F_U(2) - F_U(-1) = 0{,}97725 - (1-0{,}84134) = 0{,}97725 - 0{,}15866 = 0{,}81859.$

\medskip
\textbf{c)} $\P(X > x) = 0{,}05$ signifie $\P(U > u) = 0{,}05$, donc $u = u_{0{,}95} = 1{,}6449$.

$x = m + \sigma \cdot u = 100 + 15 \times 1{,}6449 = 124{,}67.$

% --- Exercice 2 ---
\bigskip
\begin{thmbox}[Exercice 2 -- Additivit\'e et loi du Khi-deux]
Soient $U_1, U_2, U_3$ i.i.d. $\sim \N(0;1)$. On pose $K = U_1^2 + U_2^2 + U_3^2$.
\begin{enumerate}[label=\alph*)]
\item Donner la loi de $K$.
\item Calculer $\P(K > 7{,}815)$.
\item D\'eterminer $k$ tel que $\P(K \leq k) = 0{,}95$.
\end{enumerate}
\end{thmbox}

\textbf{Correction :}

\textbf{a)} $K = \sum_{i=1}^{3} U_i^2 \sim \chi^2(3)$ (somme de 3 carr\'es de NCR ind\'ependantes).

\medskip
\textbf{b)} Table~4, ligne $\nu=3$, on cherche 7,815. On lit $\chi^2_{0{,}95}(3) = 7{,}815$ \`a la colonne $P=0{,}05$.

Donc $\P(K > 7{,}815) = 0{,}05$.

\medskip
\textbf{c)} $\P(K \leq k) = 0{,}95 \Leftrightarrow \P(K > k) = 0{,}05$.

On lit $k = \chi^2_{0{,}95}(3) = 7{,}815$.

% --- Exercice 3 ---
\bigskip
\begin{thmbox}[Exercice 3 -- Construction d'une loi de Student]
Soient $U \sim \N(0;1)$, $K_1 \sim \chi^2(4)$ et $K_2 \sim \chi^2(6)$, toutes ind\'ependantes.
\begin{enumerate}[label=\alph*)]
\item Donner la loi de $K = K_1 + K_2$.
\item Donner la loi de $T = \dfrac{U}{\sqrt{K/10}}$.
\item D\'eterminer $t$ tel que $\P(|T| > t) = 0{,}05$.
\end{enumerate}
\end{thmbox}

\textbf{Correction :}

\textbf{a)} Par additivit\'e : $K = K_1 + K_2 \sim \chi^2(4+6) = \chi^2(10)$.

\medskip
\textbf{b)} $T = \dfrac{U}{\sqrt{K/10}}$ avec $U \sim \N(0;1)$, $K \sim \chi^2(10)$ ind\'ependantes et $\nu = 10$.

Donc $T \sim t(10)$.

\medskip
\textbf{c)} La Table~5 de Student est directement bilat\'erale. Lire colonne $P=0{,}05$, ligne $\nu=10$ :

$t = T_{0{,}95}(10) = 2{,}2281.$

Les bornes sont $-2{,}2281$ et $+2{,}2281$.

% --- Exercice 4 ---
\bigskip
\begin{thmbox}[Exercice 4 -- Fisher-Snedecor]
Soient $K_1 \sim \chi^2(3)$ et $K_2 \sim \chi^2(5)$ ind\'ependantes. On pose $F = \dfrac{K_1/3}{K_2/5}$.
\begin{enumerate}[label=\alph*)]
\item Donner la loi de $F$.
\item D\'eterminer $F_{0{,}95}(3;5)$.
\item D\'eterminer $F_{0{,}05}(3;5)$ (queue gauche).
\item Trouver les bornes de l'intervalle bilat\'eral \`a $90\%$.
\end{enumerate}
\end{thmbox}

\textbf{Correction :}

\textbf{a)} $F \sim F(3;5)$ par d\'efinition.

\medskip
\textbf{b)} Table~6, colonne $\nu_1=3$, ligne $\nu_2=5$ : $F_{0{,}95}(3;5) = 5{,}41$.

\medskip
\textbf{c)} Formule d'inversion : $F_{0{,}05}(3;5) = \dfrac{1}{F_{0{,}95}(5;3)} = \dfrac{1}{9{,}01} \approx 0{,}111$.

\medskip
\textbf{d)} Intervalle bilat\'eral \`a $90\%$ signifie $P=10\%$, soit $5\%$ de chaque c\^ot\'e.
\begin{itemize}
\item Borne sup\'erieure : $F_{0{,}95}(3;5) = 5{,}41$.
\item Borne inf\'erieure : $F_{0{,}05}(3;5) = 1/F_{0{,}95}(5;3) = 1/9{,}01 \approx 0{,}111$.
\end{itemize}
$\P(0{,}111 < F \leq 5{,}41) = 90\%.$

% --- Exercice 5 ---
\bigskip
\begin{thmbox}[Exercice 5 -- Synth\`ese (trouver les fractiles)]
D\'eterminer les fractiles suivants :
\begin{enumerate}
\item $u_{0{,}144}$
\item $\chi^2_{0{,}90}(14)$
\item $T_{0{,}70}(2)$ (intervalle bilat\'eral sym\'etrique)
\item $F_{0{,}01}(3;1)$
\item $u_{0{,}601}$
\end{enumerate}
\end{thmbox}

\textbf{Correction :}

\textbf{1)} $u_{0{,}144}$ : comme $Q = 0{,}144 < 0{,}5$, on a $u_{0{,}144} = -u_{0{,}856}$. Lire dans la Table~3 \`a $Q = 0{,}856$ : $u_{0{,}856} = 1{,}0625$. Donc $u_{0{,}144} = -1{,}0625$.

\medskip
\textbf{2)} $\chi^2_{0{,}90}(14)$ : Table~4, colonne $P=0{,}1$, ligne $\nu=14$ $\Rightarrow$ $\chi^2_{0{,}90}(14) = 21{,}064$.

\medskip
\textbf{3)} $T_{0{,}70}(2)$ (bilat\'eral) : lire colonne $P=0{,}30$, ligne $\nu=2$ $\Rightarrow$ $T_{0{,}70}(2) = 1{,}3862$.

Par sym\'etrie : $T_{0{,}30}(2) = -1{,}3862$.

\medskip
\textbf{4)} $F_{0{,}01}(3;1)$ : queue gauche $\Rightarrow$ $F_{0{,}01}(3;1) = \dfrac{1}{F_{0{,}99}(1;3)}$.

Table~6, colonne $\nu_1=1$, ligne $\nu_2=3$ \`a $P=1\%$ : $F_{0{,}99}(1;3) = 34{,}12$.

$F_{0{,}01}(3;1) = 1/34{,}12 \approx 0{,}0293$.

\medskip
\textbf{5)} $u_{0{,}601}$ : comme $Q = 0{,}601 > 0{,}5$, lecture directe Table~3 \`a $Q = 0{,}601$ $\Rightarrow$ $u_{0{,}601} = 0{,}2559$.

% ============================================================
% VIII. PIEGES CLASSIQUES
% ============================================================
\newpage
\section{Pi\`eges classiques}

\begin{alertbox}[Pi\`ege 1 -- Confondre $\sigma$ et $\sigma^2$]
$\N(m;\sigma)$ : le deuxi\`eme param\`etre est l'\textbf{\'ecart-type}, pas la variance. Certains manuels utilisent $\N(m;\sigma^2)$. V\'erifier syst\'ematiquement la convention de l'\'enonc\'e.
\end{alertbox}

\begin{alertbox}[Pi\`ege 2 -- Confondre $f(-u) = f(u)$ et $F(-u) = 1-F(u)$]
La \textbf{densit\'e} est sym\'etrique : $f(-u) = f(u)$.\\
La \textbf{fonction de r\'epartition} v\'erifie : $F(-u) = 1 - F(u)$.\\
Ne \textbf{jamais} \'ecrire $F(-u) = F(u)$.
\end{alertbox}

\begin{alertbox}[Pi\`ege 3 -- Oublier de diviser $P$ par 2 en bilat\'eral (Normale et Khi-deux)]
Pour les tables de la loi normale et du Khi-deux, un intervalle bilat\'eral sym\'etrique \`a risque $P$ se lit \`a $P/2$. Seule la table de Student est \textbf{d\'ej\`a bilat\'erale}.
\end{alertbox}

\begin{alertbox}[Pi\`ege 4 -- Table de Student : unilat\'eral vs bilat\'eral]
La Table~5 est pr\'evue pour le bilat\'eral. Pour un test \textbf{unilat\'eral} \`a risque $P$, il faut lire \`a la colonne $2P$.
\end{alertbox}

\begin{alertbox}[Pi\`ege 5 -- Inverser $\nu_1$ et $\nu_2$ dans Fisher]
$F(\nu_1;\nu_2) \neq F(\nu_2;\nu_1)$ en g\'en\'eral. Toujours respecter : $\nu_1$ = ddl du num\'erateur, $\nu_2$ = ddl du d\'enominateur. Dans la table : $\nu_1$ en \textbf{colonne}, $\nu_2$ en \textbf{ligne}.
\end{alertbox}

\begin{alertbox}[Pi\`ege 6 -- Oublier l'hypoth\`ese d'ind\'ependance]
Les constructions $\chi^2$, Student et Fisher \textbf{exigent} l'ind\'ependance entre les variables al\'eatoires. Sans ind\'ependance, les lois ne sont pas valides.
\end{alertbox}

\begin{alertbox}[Pi\`ege 7 -- Additivit\'e des \'ecarts-types]
Ce sont les \textbf{variances} qui s'additionnent, pas les \'ecarts-types.
$$\sigma_{X+Y} = \sqrt{\sigma_X^2 + \sigma_Y^2} \;\neq\; \sigma_X + \sigma_Y.$$
\end{alertbox}

\end{document}